\section{Related Work}

\subsection{Transformer and One-Stream Visual Tracking}
Recent RGB trackers commonly formulate single-object tracking as template-search interaction. Early transformer trackers used attention to connect the initial target template with the current search region and to improve discriminative relation modeling~\cite{chen2021transt,yan2021stark}. Subsequent one-stream trackers further reduced the separation between feature extraction and relation modeling. MixFormer uses mixed attention for simultaneous feature extraction and target-information integration, while \baseline{} unifies feature learning and relation modeling by passing template and search tokens through a shared stream~\cite{cui2022mixformer,ye2022ostrack}.

Other recent trackers explore complementary formulations. SeqTrack treats tracking as sequence generation, ROMTrack models inherent and hybrid template features, and HIPTrack uses historical prompts~\cite{chen2023seqtrack,cai2023romtrack,cai2024hiptrack}.

ODTrack propagates dense temporal tokens online, LoRAT adapts large vision transformers with low-rank fine tuning, and AQATrack uses autoregressive spatio-temporal queries~\cite{zheng2024odtrack,lin2024lorat,xie2024aqatrack}.

These methods show that modern tracking performance often depends on how target information is represented, propagated, and fused with the search region. They also provide the architectural context for this paper. The present work does not propose a new transformer tracker, a new generative tracking head, or a new large-backbone adaptation strategy. Instead, it keeps the \baseline{} template-search formulation and studies whether the already extracted search representation can be adapted under degraded observations.

\subsection{Temporal Modeling and Target Adaptation}
Temporal information is a major source of robustness in visual tracking because target appearance, background clutter, and localization uncertainty change over time. STARK introduced spatio-temporal transformer tracking, while HIPTrack and ODTrack show two recent directions for using tracking history through historical prompts or dense temporal token propagation~\cite{yan2021stark,cai2024hiptrack,zheng2024odtrack}. SeqTrack and AQATrack also connect tracking with autoregressive or spatio-temporal query mechanisms, replacing part of the traditional hand-designed tracking pipeline with sequence or query prediction~\cite{chen2023seqtrack,xie2024aqatrack}. These designs are relevant because they improve target adaptation without requiring the present paper's restoration-guided state-space feature block.

The distinction is important for the proposed method. Temporal prompts, autoregressive queries, and online token propagation aim to exploit historical target information or simplify tracking prediction. They do not by themselves establish whether a degraded search feature can be made more consistent with its clean counterpart after search-token extraction. The final \method{} configuration therefore avoids memory update, response fusion, target-region consistency, and target-versus-distractor margin loss, because the retained evidence supports a narrower global clean--degraded feature-consistency setting rather than a broader temporal-adaptation design.

\subsection{State-Space Models for Visual Representation}
State-space sequence modeling has recently become an alternative mechanism for long-range dependency modeling. Mamba introduces selective state spaces with input-dependent parameters and hardware-aware computation for efficient sequence modeling~\cite{gu2024mamba}. Vision Mamba adapts bidirectional Mamba blocks to visual representation learning, and VMamba develops visual state-space blocks with two-dimensional selective scanning for vision backbones~\cite{zhu2024visionmamba,liu2024vmamba}. These works are not tracking methods, but they are important because they show how state-space models can be reorganized for visual tokens rather than only for one-dimensional language or signal sequences.

Restoration-oriented Mamba models adapt this family in a different direction. MambaIR addresses low-level restoration with residual state-space blocks, local enhancement, and channel attention, while MambaIRv2 develops attentive state-space restoration and semantic-guided neighboring to reduce limitations of causal scanning in image restoration~\cite{guo2024mambair,guo2025mambairv2}. These papers support the technical motivation for using state-space processing under degraded visual evidence. However, their objectives are image restoration objectives, not template-search localization objectives. A tracker needs target-discriminative features and accurate box prediction, so an external restoration network or full image-reconstruction objective cannot be assumed to improve tracking.

\subsection{Mamba-Based Visual Tracking}
Mamba-based tracking has already been explored, so the present paper does not claim novelty from simply using Mamba in tracking. MambaLCT uses a Context Mamba module to aggregate long-term target variation cues over video frames~\cite{li2024mambalct}. Mamba-FETrack and Mamba-FETrack V2 use Mamba-style backbones and fusion mechanisms for RGB-event tracking, where event data provide complementary sensing under difficult imaging conditions~\cite{huang2024mambafetrack,wang2025mambafetrackv2}. MambaEVT studies event-stream tracking with Vision Mamba and Memory Mamba for dynamic template generation~\cite{wang2024mambaevt}. MambaNUT applies a pure Mamba one-stream design with adaptive curriculum learning for nighttime UAV tracking~\cite{wu2024mambanut}.

These works demonstrate several uses of state-space modeling in tracking: temporal context aggregation, multimodal RGB-event fusion, event-only dynamic template generation, efficient backbone design, and nighttime template-search learning. They are therefore direct novelty boundaries for this paper. The proposed \method{} is not a complete Mamba backbone, not an RGB-event or event-only tracker, not a temporal-memory module, and not a nighttime-specific curriculum tracker. Its placement is more limited: it is inserted after search-token extraction in \baseline{} to adapt degradation-sensitive search representations while retaining the frozen backbone and training only \texttt{rgssb.*} and \texttt{box\_head.*}.

\subsection{Tracking Under Degraded Observations and Present-Work Positioning}
Tracking under degraded observations has been studied through benchmark construction, adverse-condition tracking, and degradation-invariant learning. LLOT explicitly targets low-light object tracking and reports a benchmark intended to expose performance gaps under low illumination~\cite{zhong2024llot}. MambaNUT further confirms that nighttime UAV tracking can be handled as a specific low-light/adverse-condition problem without using a separate restoration module~\cite{wu2024mambanut}. The local InvTrack evidence is also central to the boundary of this paper because it studies degradation-invariant single-object tracking with clean/degraded branches, feature consistency, and response-map fusion~\cite{invtrack_local}. Thus, the manuscript does not claim that low-light tracking, degradation-aware training, feature consistency, or response-map fusion is new.

The present contribution is instead a controlled test of a narrower intersection. It differs from replacing \baseline{} with a complete Mamba tracker, from multimodal Mamba trackers that rely on event data, from temporal-memory Mamba trackers, from nighttime-specific Mamba trackers, from image restoration used as external preprocessing, and from full restoration networks optimized for image quality. The proposed \method{} adapts search features inside the tracker after search-token extraction and uses global clean--degraded feature consistency with weight $\lambda=0.02$. This positioning supports only a limited claim: restoration-guided state-space adaptation can be evaluated as a feature-level modification to \baseline{} under selected RGB degradations. It does not establish state-of-the-art tracking performance, universal degradation robustness, or uniform improvement across all benchmarks.
